\section{What release 1.4 adds}
\label{sec:release14}

Release 1.4 removes the single-process ceiling of \cref{sec:limitations} and connects a model
to arriving data, to continuous fields and to the browser. This section gives, for each
addition, what was measured and where it falls short; \cref{tab:release14} collects both, and
\cref{app:release14} has the designs, the measurement conditions and the files behind every
number. These measurements were made on the laptop of \cref{sec:performance} while it was
shared with other work and without the fixed frequency of \cref{sec:playbooks}; the
conclusions that matter (byte-identical outputs, order-of-magnitude speed-ups, shortfalls by
a factor) do not turn on the last digit of a timing.

\paragraph{Per-agent random streams and split runs.}
\label{sec:split}
A model can give every agent random generators of its own, seeded from the scenario's seed
and the agent's key, so that what an agent draws no longer depends on the activation order or
on the population (the flag is off by default). On that basis one run can be split across
workers: a grid into blocks with halos, a network by breadth-first banding, a continuous space
by recursive bisection, with ghost agents exchanged before a phase that reads across a border
and movers migrating after a phase that moves them, over local, MPI, Dask or Ray workers. On
three reference models (a grid epidemic, a network SIR, a continuous SIR) every reporter and
every agent's state is byte-identical to the single-process run at 1, 2, 4 and 16 partitions,
under MPI with 2 and 4 ranks and on two Dask workers (\cref{app:split}). The aim of 70\%
weak-scaling efficiency from 1 to 64 workers could not be measured on one laptop and is not
met on it: 62\% at 2 workers, 20\% at 8 and 9\% at 16 (\cref{tab:scaling}), where eight to
sixteen copies of the same work with no communication at all already run at 23 to 32\% of
one, because beyond the four performance cores the copies land on efficiency cores; the
protocol keeps 53 to 91\% of what the machine gives, and 79 to 109\% with 40,000 people per
worker. The Ray backend is written but untested, as is MPI on macOS and Windows.

\paragraph{A GPU backend.}
\label{sec:gpu}
Column phases, cohort tables and the cell-space operators run on CuPy or JAX with a
\code{numpy} fallback, from counter-based (Philox) random streams that draw the same numbers
on every backend. At ten million agents the grid epidemic's 60 steps take 0.61~s on an
RTX~5050 laptop GPU against 21~s for the CPU column variant, 34.4 times faster (three rounds
each), with the same infected series and the same digest of the final state on \code{numpy}
and on CuPy; building ten million agent objects takes 22 to 42~s on either side, so build plus
steps gains only 1.5 to 1.8 times (\cref{app:gpu}). Below tens of thousands of agents the GPU
is slower than the CPU (0.17 to 0.21~s against 0.11~s at the playbook's own 20,000 agents), a
cohort table of under 200 rows is 200 times slower, and Apple Metal was not measured.

\paragraph{Differentiable models.}
\label{sec:diff}
Column phases traced by JAX give gradients of a model's outputs with respect to its
parameters, exact through normal and uniform draws and approximate through Bernoulli choices
(straight-through, Gumbel-softmax or score-function estimators, each warning that it is one).
Gradient calibration recovers the Bass diffusion model's rates (0.0305 and 0.379 for 0.03 and
0.38) and the Bouchaud--Mézard model's (1.448 and 0.937 for 1.5 and 1.0; 1.54 and 1.09 with
four replicates) from runs of the agent-based zoo models, and pathwise gradients agree with
finite differences to a relative error below $10^{-6}$ on a one-parameter test model and
below $10^{-4}$ on the Bouchaud--Mézard twin (\cref{app:diff}); the straight-through estimator
is biased and fails that check, by design.

\paragraph{Live data and data assimilation.}
\label{sec:assimilation}
Observations arrive from a replayed table, a growing CSV file, an HTTP endpoint or a message
broker; an ensemble of replicates advances in lock step, a bootstrap particle filter
\citep{gordon1993novel} weights and resamples it, and an ensemble Kalman filter
\citep{evensen2003ensemble} adjusts continuous parameters, both deterministic for a seed. On a
synthetic town of 1,000 people the filtered nowcast of the unobserved number infected has a
root-mean-square error of 5.07 people against 69.2 for the unfiltered ensemble over 50 days,
and the Kalman filter recovers the contact rate as 0.451 for a truth of 0.45
(\code{paper/data/assimilation.json}, \cref{app:assimilation}). Message brokers were tested
against fake clients only, and stepping stays in one process.

\paragraph{Fields and coupling.}
\label{sec:couple}
A field is a property layer with a solver (diffusion, upwind advection and decay, in explicit
sub-steps that refuse an unstable time step, or backward Euler), coupled to differential
equations and stock-flow components under a schedule that is checked when the model is built.
Closed and periodic boundaries conserve mass to a relative $10^{-12}$, a spreading Gaussian
matches the heat kernel to about $10^{-6}$, and two zoo models use it: foragers on a diffusing
nutrient (a population within 1\% of the carrying capacity of its energy balance, the
nutrient budget closed to round-off) and an epidemic coupled to hospital capacity (the
final-size relation within five standard deviations, the hospital's exact matrix-exponential
solution within $10^{-6}$) (\cref{app:couple}). Fields need orthogonal grids with uniform
spacing, and upwind advection smears sharp fronts.

\paragraph{Emulators, active learning and history matching.}
\label{sec:emulate}
A Gaussian-process or neural emulator trained on an experiment's runs, active learning that
picks the next runs by integrated variance reduction, and history matching in waves
\citep{vernon2010galaxy, andrianakis2015bayesian}. On the age-structured SEIR model four waves
of 20 runs leave 17.1, 4.4, 3.1 and 2.5\% of the parameter box, with the true parameters
inside after every wave (\cref{app:emulate}). Active learning needs half the runs of a Latin
hypercube for the same emulator error on the Bass diffusion model at most budgets (factors of
1.3 to 2.7) but not on the Branin (1.1 to 1.4) or Ishigami (0.8 to 1.4) functions, where a
well-spread design is already close to the best that a variance criterion blind to the
outputs can do.

\paragraph{Synthetic populations and models on real data.}
\label{sec:population}
Iterative proportional fitting, integerisation, households, places within a distance limit and
a contact network with household, school and work layers build a starting population from
aggregate tables; no personal data is bundled, the marginals of the tests are made up, and
three small real samples (OpenStreetMap streets, ESA WorldCover land cover, census tracts)
ship with their licences. About 20,000 synthetic people reproduce their totals and age counts
exactly in every tract and their sex counts within 1\%, the SEIR model started from them
follows the final-size relation within 0.05, and three zoo models run on the real samples,
each checked against a known answer (\cref{app:population}).

\paragraph{\modelly{} in the browser.}
\label{sec:browser}
\cmd{modelly export-web} writes a static site that runs a model with Pyodide in a browser tab,
with Studio's parameter panel, a space view and a reporter plot. The 14 bundled examples of
release 1.3 and 12 zoo models each ran from their export in headless Chromium with the
pure-Python wheel and with the compiled one (52 of 52 runs); the five zoo models added later
had not run in a browser when this was written (\cref{app:browser}). A cold first visit is
ready after 8.0~s on a 50~Mbit/s line for a model that needs no \code{scipy}, 12.3~s for one
that does, and 21.4~s on a Fast~4G profile, so the aim of under ten seconds holds only on the
faster line without \code{scipy}.

\paragraph{Faster events and data collection, columns on disk.}
\label{sec:perf2}
The event list's heap is a fifth C module with a pure-Python twin, the data collector reads
column attributes in C, and a table's columns can live in memory-mapped files. Against 1.3.3
the event-heavy wolf-sheep example runs at 0.95 of the time, a schedule of 600 events a step
at 0.70 and a collection over 20,000 column agents at 0.47, with identical results; the
playbooks are unchanged (0.986 to 1.004), and a cohort table of 100 million rows ran within a
16~GB memory limit at a peak of 3.6~GB (\cref{app:perf2}). The predator-prey playbook, which
schedules no events, is not faster; profile-guided optimisation measured no gain and is not
shipped; and 100 million agent objects do not fit in 16~GB.

\begin{table}[htbp]
  \centering
  \footnotesize
  \caption{Release 1.4: what was measured and where it falls short. Details and sources in
  \cref{app:release14}.}
  \label{tab:release14}
  \begin{tabularx}{\linewidth}{>{\raggedright\arraybackslash}p{0.17\linewidth}>{\raggedright\arraybackslash}X>{\raggedright\arraybackslash}X}
    \toprule
    Capability & Measured & Falls short \\
    \midrule
    Split runs & Byte-identical to one process at 1, 2, 4 and 16 partitions on three models;
    MPI with 2 and 4 ranks; Dask with 2 workers & Weak scaling 62, 20 and 9\% at 2, 8 and 16
    workers with 10,000 people per worker (67 and 28\% at 2 and 8 with 40,000) on one laptop;
    64 workers unmeasured; Ray, and MPI on macOS and Windows, untested \\
    \addlinespace
    GPU backend & Steps 34.4 times faster at 10 million agents, identical integer state &
    Construction not faster (1.5 to 1.8 times in total); slower below tens of thousands of
    agents; cohort tables 200 times slower; Metal untested \\ \addlinespace
    Differentiable models & Two zoo models' parameters recovered within 2\% and 7\%; pathwise
    gradients within $10^{-6}$ and $10^{-4}$ of finite differences & Gradients through discrete
    choices are approximations; the straight-through estimator is biased \\ \addlinespace
    Data assimilation & Nowcast error 5.07 against 69.2 unfiltered; contact rate 0.451 for
    0.45; bit-identical for a seed & Brokers tested with fake clients; replicates step in one
    process \\ \addlinespace
    Fields and coupling & Mass conserved to $10^{-12}$; heat kernel to about $10^{-6}$; two
    zoo models with closed budgets & Orthogonal grids with uniform spacing only; upwind
    advection smears fronts \\ \addlinespace
    Emulators & Parameter box reduced from 17.1 to 2.5\% in four waves, truth inside & Half the
    runs of a hypercube only on the zoo model (factors 1.3 to 2.7); Branin 1.1 to 1.4, Ishigami
    0.8 to 1.4 \\ \addlinespace
    Synthetic populations & Totals and ages exact, sexes within 1\%; the SEIR model's final
    size within 0.05; three real-data zoo models & Made-up marginals; no real joint
    distribution to test against \\ \addlinespace
    Browser & 52 of 52 runs of 26 models, two wheels & Five zoo models not yet run; first
    load 8.0, 12.3 and 21.4~s by line and model, under 10~s only without \code{scipy} at
    50~Mbit/s \\ \addlinespace
    Events, collection, columns on disk & Wolf-sheep 0.95, events 0.70, collection 0.47 of
    1.3.3's time; playbooks 0.986 to 1.004; 100 million rows within 16~GB & Predator-prey
    unchanged; profile-guided optimisation no gain; 100 million agent objects need over 22~GB
    \\ \addlinespace
    Platforms & Every check on Linux x86-64 & Three-platform determinism run and the import
    check of the macOS, Windows and ARM wheels not repeated for 1.4.0; free-threaded build of
    the event heap untested \\
    \bottomrule
  \end{tabularx}
\end{table}
